\documentclass[11pt,a4paper]{article}
\usepackage[margin=2.5cm]{geometry}
\usepackage{amsmath,amssymb,amsthm}
\usepackage{booktabs}
\usepackage{cite}
\usepackage[hidelinks]{hyperref}

\newcommand{\R}{\mathbb{R}}
\newcommand{\C}{\mathbb{C}}
\newcommand{\A}{\mathcal{A}}
\newcommand{\B}{\mathcal{B}}
\newcommand{\G}{\mathcal{G}}
\newcommand{\J}{\mathcal{J}}
\newcommand{\Rt}{\widetilde R}
\newcommand{\todo}[1]{\textbf{[TODO: #1]}}

\newtheorem{theorem}{Theorem}[section]
\newtheorem{lemma}[theorem]{Lemma}
\newtheorem{proposition}[theorem]{Proposition}
\newtheorem{assumption}[theorem]{Assumption}
\theoremstyle{remark}
\newtheorem{remark}[theorem]{Remark}

\title{Preconditioning for quantum linear solvers\\ in PDE-constrained optimization}
\author{Martin Stoll \and John W. Pearson}
\date{Working draft, November 2026}

\begin{document}
\maketitle

\begin{abstract}
Quantum linear system solvers cost a number of operations proportional to an
\emph{effective} condition number. This number is the ratio of the
normalization with which the matrix is embedded in a unitary circuit (its
block encoding) to its smallest singular value. Preconditioning is
therefore the main lever that numerical linear algebra has on quantum
solvers. The rules differ from the classical ones, however:
\begin{itemize}
\item clustering of eigenvalues does not help;
\item preconditioners must be applied by multiplication, not by inner
  solves;
\item every factor of a product must be well scaled.
\end{itemize}
We develop preconditioners under these rules for the saddle-point systems
of PDE-constrained optimization. For distributed control of the Poisson
equation, we combine a Sch\"oberl--Zulehner block preconditioner with an
additive multilevel (BPX) factor. We apply it on both sides of the KKT
matrix and factor the result as $\G^*\J\G$, where $\J$ is a single
two-qubit reflection and $\G$ is well scaled. Building on the factorized
multilevel block encodings of Deiml and Peterseim, which we formulate
through quadrature points so that they cover $Q_k$ and $P_k$ elements, $\G$
has a block encoding whose normalization grows only like $\sqrt{\log h^{-1}}$. The effective condition number is then bounded
independently of the mesh size. With level weights that we introduce, it
is also bounded in the regularization parameter, up to a constant fixed by
the mass matrix. For the time-periodic heat equation, a quantum Fourier
transform in time decouples the problem into frequency-wise saddle-point
problems of the same form. We prove a $\sqrt3$ bound uniform in all
parameters, and the coupling matrix becomes a scaled reflection.
Numerical experiments in one, two and three dimensions confirm the
bounds. We also discuss when an asymptotic quantum advantage can and
cannot be expected.
\end{abstract}

\section{Introduction}\label{sec:intro}

Quantum linear system algorithms (QLSAs) prepare a quantum state proportional
to the solution $x$ of $Ax=b$. They use a number of calls to a circuit
encoding $A$ that scales like $\kappa\log(1/\varepsilon)$, where $\varepsilon$ is the accuracy of
the prepared solution state \cite{CKS17,GSLW19,Costa22}. The
size $n$ of the system enters only through the circuit cost, which is
typically polylogarithmic in $n$ for structured problems. Discretized PDEs are
the obvious application, but the condition number of a discretized
differential operator grows with the resolution, e.g. like $h^{-2}$ for
second-order elliptic problems. Without preconditioning, the dependence on
$\kappa$ cancels the gain in $n$ \cite{MP16,LMS22}.

Preconditioning for QLSAs has been studied:
\begin{itemize}
\item by sparse approximate inverses \cite{CJS13};
\item by ``fast inversion'' of a dominant part of the matrix \cite{TAWL21};
\item by wavelet bases \cite{BNWA23};
\item for elliptic finite element discretizations, by Deiml and Peterseim
  \cite{DP26}. Their work is the basis of our Section~\ref{sec:poisson}.
  \begin{itemize}
  \item They write the stiffness matrix as $C_L^T(D_A\otimes I)C_L$, where $C_L$ maps
    nodal values to gradients in an $L^2$-orthonormal basis of discontinuous
    $Q_1$ functions, and $D_A$ holds the (possibly variable) diffusion
    coefficient.
  \item They apply the BPX preconditioner as a generating system $F$, with
    $FF^T$ the preconditioner, to this factor.
  \item A commutative diagram (gradient on the coarse level, followed by an
    inclusion with orthonormal columns) gives a block encoding of the
    preconditioned factor whose normalization grows only linearly in the
    number of levels.
  \item They introduce the effective condition number (subnormalization
    times condition number), and they ran small instances on quantum
    hardware.
  \end{itemize}
  Related constructions are the finite-difference analogue \cite{JLMY25} and
  the extension to reaction--diffusion problems $-\Delta u+cu$ with a bound
  uniform in $c$ \cite{YYZ26};
\item empirically, for CFD matrices, comparing several preconditioners and
  ways of encoding them \cite{LS25};
\item for the Newton systems of quantum interior point methods
  \cite{QIPMreview}, which become increasingly ill conditioned along the
  central path.
\end{itemize}
These works make clear that preconditioning for a quantum solver is a
different problem from classical preconditioning
(Section~\ref{sec:prelim}). Only a bound on the condition number counts,
not the distribution of the eigenvalues. The preconditioner must be
realized as a circuit that \emph{multiplies}. And the normalizations of
all factors enter the cost.

\paragraph{Gap.} The saddle-point systems of PDE-constrained optimization
have not been considered from this angle. Classically, they are the
prototype for parameter-robust block preconditioners
\cite{RDW10,SZ07,PW12}, which are robust in the mesh size $h$ and the
regularization parameter $\beta$. Quantum approaches to PDE-constrained
optimization have so far been variational heuristics.

\paragraph{Contributions.}
\begin{enumerate}
\item We formulate design principles for preconditioners of QLSAs in matrix
  terms, centred on the effective condition number $\kappa_\alpha$
  (Section~\ref{sec:prelim}).
\item For distributed control of the Poisson equation
  (Section~\ref{sec:poisson}), we combine the Sch\"oberl--Zulehner block
  preconditioner with a diagonally scaled BPX factor $R$, applied on both
  sides.
  \begin{itemize}
  \item The resulting matrix factors as $\G^*\J\G$. Here $\J$ is a single
    two-qubit reflection, and $\G$ maps multilevel coefficients to values
    and gradients at quadrature points.
  \item The block encoding of $\G$, with normalization $O(\sqrt{\log h^{-1}})$, builds
    on the commuting diagram of Deiml and Peterseim \cite{DP26}, extended to
    the mass term in \cite{YYZ26}. We formulate it through quadrature
    points, with an \emph{element-local isometry} (Lemma~\ref{lem:commute}),
    so that it covers $Q_k$ and $P_k$ elements.
  \item The effective condition number is bounded independently of $h$
    (Theorem~\ref{thm:main}). It holds for $Q_k$ and $P_k$ elements on
    nested meshes.
  \end{itemize}
\item We adapt robust multilevel preconditioning for $I-\sqrt\beta\,\Delta$
  \cite{GO95,BPV00} to the Jacobi-scaled factored form that the block
  encoding needs. A family of normalized level weights makes the spectral
  interval $[c_0,c_1]$ itself independent of $h$ and $\beta$
  (Theorem~\ref{thm:robust}, Section~\ref{sec:weights}), with a
  self-contained proof and explicit constants.
\item For the time-periodic heat equation (Section~\ref{sec:heat}), the QFT
  in time reduces the problem to frequency-wise problems of the Poisson
  type, with effective parameter $\beta_k=\beta/r_k^2$.
  \begin{itemize}
  \item We prove a bound of $\sqrt3$ for the block preconditioner, uniform
    in $h$, $\tau$, $\beta$ and the frequency (Theorem~\ref{thm:heatsz}).
  \item The coupling matrix becomes a scaled reflection of norm at most
    one.
  \item Neither the number of time steps nor the step size enters the
    effective condition number.
  \end{itemize}
\item We verify all bounds and factorizations numerically in one, two and
  three dimensions (Section~\ref{sec:num}). We give an honest discussion of
  complexity: an asymptotic advantage over optimal classical solvers
  requires $d=3$ for the elliptic problem. For the time-periodic problem it
  requires data that are rough in time (Section~\ref{sec:discussion}).
\end{enumerate}

\section{Quantum linear algebra in matrix terms}\label{sec:prelim}

We describe the computational model in the language of numerical linear
algebra (see \cite{GSLW19,Lin22} for details).

\paragraph{States and outputs.} A unit vector $v\in\C^{2^k}$ is stored in $k$
qubits. Its entries cannot be read out. Repeated measurement gives
estimates of quantities $v^*Ov$ for suitable Hermitian $O$, with additive error
$\eta$. This costs $O(\eta^{-2})$ repetitions, or $O(\eta^{-1})$ with amplitude
estimation \cite{BHMT02}. The output of a quantum solver is therefore a
small number of scalar \emph{quantities of interest} (QoIs), such as
$\langle w,y\rangle_M$ or the value of the objective, never the vector $x$.

\paragraph{Block encodings.} A matrix $A$ with $\|A\|\le\alpha$ is \emph{block encoded}
with \emph{subnormalization} $\alpha$ if a unitary circuit $U$ satisfies
\[
U=\begin{bmatrix}A/\alpha & *\\ * & *\end{bmatrix}.
\]
Every block of a unitary matrix has norm at most one, so $\alpha\ge\|A\|$ for
\emph{every} block encoding. We call $\alpha=\|A\|$ the \emph{ideal
subnormalization}. It is always attainable in principle, by a unitary
dilation, but that dilation is a dense matrix without an efficient
circuit. Efficient circuits are built from structure, and they usually pay
more than $\|A\|$. How much more depends on the construction; examples are in
the supplement (\texttt{nla\_for\_quantum.tex}). Applying $U$ to $[v;0]$ and keeping the first block
(postselection) yields $Av/\|Av\|$. This succeeds with probability
$\|Av\|^2/(\alpha\|v\|)^2$.

A circuit is a product of unitary matrices (gates), each acting on a few
of the tensor factors $\C^2$. For small problems, the whole product can be
formed as an explicit matrix and checked. As an example, let $S$ be the
cyclic shift (a permutation matrix) and $H=\frac1{\sqrt2}\begin{bmatrix}1&1\\1&-1\end{bmatrix}$. Then
\[
U=(H\otimes I)\begin{bmatrix}I&0\\0&-S\end{bmatrix}(H\otimes I)
=\frac12\begin{bmatrix}I-S & I+S\\ I+S & I-S\end{bmatrix}
\]
is unitary, and it block encodes the difference matrix $I-S$ with $\alpha=2$.
This is the simplest instance of a linear combination of unitaries (LCU):
$H$ prepares the weights, and the middle factor selects $I$ or $-S$ according
to one extra qubit.
Rectangular matrices are padded with zeros to square matrices of
power-of-two size. Block encodings combine as follows (subnormalizations in
brackets):
\begin{itemize}
\item products $AB$ $[\alpha_A\alpha_B]$;
\item Kronecker products $A\otimes B$ $[\alpha_A\alpha_B]$;
\item linear combinations $\sum_jc_jA_j$ $[\sum_j|c_j|\alpha_j]$;
\item stacking $[A_1;\dots;A_k]$ $[(\sum_j\alpha_j^2)^{1/2}]$: prepare an index
  register with amplitudes $\propto\alpha_j$ and apply $A_j$ controlled on
  index $j$;
\item concatenation $[A_1,\dots,A_k]$ $[(\sum_j\alpha_j^2)^{1/2}\le\sqrt k\max_j\alpha_j]$;
\item sparse matrices with $s_r$ nonzeros per row and $s_c$ per column,
  whose entries are computable by a circuit $[\sqrt{s_rs_c}\max_{ij}|a_{ij}|]$
  \cite{GSLW19}.
\end{itemize}
Products need a fresh ancilla per factor; reusing one ancilla brings
discarded components back into the leading block. Explicit unitary
realizations of all rules, with derivations of the subnormalizations and
numerical checks, are given in the supplement \emph{Block encodings for
numerical linear algebraists} (\texttt{nla\_for\_quantum.tex}).
In a product, the subnormalization can be far larger than the norm:
$\alpha_A\alpha_B\ge\|A\|\|B\|\gg\|AB\|$ is possible.

\paragraph{Solving.} Let $\B$ be Hermitian and block encoded with subnormalization
$\alpha$, and let its nonzero eigenvalues satisfy $|\lambda|\ge\sigma_{\min}$. A quantum solver
for $\B z=b$ outputs a unit vector $\tilde z$ (a quantum state). Its
\emph{accuracy} $\varepsilon$ is measured in the Euclidean norm against the normalized
minimum-norm solution:
\begin{equation}\label{eq:eps}
\Big\|\tilde z-\frac{\B^+b}{\|\B^+b\|}\Big\|\le\varepsilon .
\end{equation}
Throughout, $\varepsilon$ denotes this solver accuracy. The additive error of a
measured scalar is $\eta$ (above). The target accuracy of a quantity of
interest is $\mathrm{tol}$ (Section~\ref{sec:discussion}). Quantum
singular value transformation (QSVT) \cite{GSLW19} applies an odd polynomial
$p(\B/\alpha)$ of degree $D$ using $D$ calls to the block encoding.

\emph{Why odd.} QSVT alternates calls to $U$ and $U^*$. It acts on the
singular values $|\lambda_i|$ of $\B=Q\Lambda Q^*$, and $D$ calls give a polynomial
$p$ of degree $D$ and parity $D$ (as in $(\B^*\B)^k$ versus $\B(\B^*\B)^k$). The signs
of the eigenvalues enter only through their $D$-th power:
\[
Q\operatorname{diag}\big(\operatorname{sign}(\lambda_i)^D\big)\,p(|\Lambda|)\,Q^* .
\]
For odd $D$ this is $Qp(\Lambda)Q^*=p(\B)$. For even $D$ the signs are lost, so an
even $p$ cannot approximate $1/\lambda$ on an indefinite spectrum.

\emph{Solving means evaluating a polynomial.} A QSVT solver does not iterate.
It evaluates the matrix function $p(\B/\alpha)$ for a fixed odd polynomial with
\[
p(t)\approx\frac{\delta}{t}\ \text{ on }[\delta,1],\qquad |p|\le1\ \text{ on }[-1,1],\qquad
\delta=\frac{\sigma_{\min}}{\alpha}.
\]
The bound $|p|\le1$ is needed because $p(\B/\alpha)$ must again be a block of a
unitary. Since $p(0)=0$, the result is the scaled pseudoinverse
$\delta(\B/\alpha)^+=\sigma_{\min}\B^+$. The cost of a solve has two factors.
\begin{itemize}
\item \emph{Cost per run (the degree).} In classical terms this is Chebyshev semi-iteration whose degree is
  fixed in advance, namely $O(\kappa_\alpha\log(\kappa_\alpha/\varepsilon))$. The $\log(\kappa_\alpha/\varepsilon)$
  arises because $p$ must approximate $\delta/t$ to about $\varepsilon/\kappa_\alpha$ for the
  normalized output to satisfy~\eqref{eq:eps}. The polynomial
  depends only on the interval, not on $b$ or the actual spectrum. Unlike CG
  or MINRES, it cannot adapt.
\item \emph{Number of runs (the success probability).} Postselecting $\sigma_{\min}\B^+b$ succeeds with probability
  \[
  a^2=\frac{\|\sigma_{\min}\B^+b\|^2}{\|b\|^2}\ \ge\ \frac1{\kappa_\alpha^2}.
  \]
  Simply repeating the run until it succeeds would cost about $1/a^2$ runs.
  \emph{Amplitude amplification} \cite{BHMT02} needs only about $\pi/(4a)\le\kappa_\alpha$
  rounds. It uses two reflections: one flips the sign of the wanted
  block, the other reflects about the output vector of the circuit (one
  application of the circuit and its adjoint). Their product is a
  rotation by $2\arcsin a$ in the plane spanned by the wanted and the
  unwanted part, so $k$ rounds turn the success amplitude $a$ into
  $\sin((2k+1)\arcsin a)$, which is close to one for $k\approx\pi/(4a)$.
  In polynomial terms this is again QSVT, now of the scalar $a$ with the
  odd Chebyshev polynomial $T_{2k+1}$. Details are in the supplement.
\end{itemize}
Multiplying the two, plain QSVT inversion costs $O(\kappa_\alpha^2\log(\kappa_\alpha/\varepsilon))$ calls.
Variable-time amplitude amplification \cite{CKS17} or the discrete adiabatic
method \cite{Costa22} achieve the optimal
\begin{equation}\label{eq:cost}
O\big(\kappa_\alpha\log(1/\varepsilon)\big),\qquad \kappa_\alpha:=\frac{\alpha}{\sigma_{\min}}\ \ge\ \kappa(\B).
\end{equation}
In both cases the cost is governed by $\kappa_\alpha$, the \emph{effective condition
number} in the terminology of Deiml and Peterseim \cite[Def.~4.1]{DP26}.

\paragraph{What preconditioning means.} If solving means evaluating a
polynomial, then preconditioning means \emph{choosing the matrix in which the
polynomial is evaluated}, so that a polynomial of low degree suffices.

\emph{Why this is the only lever.} The polynomial has to approximate $\delta/t$
on $[\delta,1]$ while staying bounded by one on $[-1,1]$, where $\delta=\sigma_{\min}/\alpha=1/\kappa_\alpha$.
This needs degree about $\kappa_\alpha\log(1/\varepsilon)$, and every degree is one call to
the block encoding. Solving is cheap only if $\delta$ is not small: the
\emph{scaled} spectrum of the encoded matrix must stay away from zero.

\emph{How it is done.} Let $\A x=f$ be the system and let $\Rt$ be a factor of
the preconditioner, $\Rt\Rt^T\approx P^{-1}$. Instead of evaluating a polynomial in
$\A/\alpha_\A$, we evaluate one in the preconditioned matrix $\B=\Rt^T\A\Rt$ and map
back:
\[
\A^{-1}=\Rt\,\B^+\Rt^T\ \approx\ \frac1{\sigma_{\min}(\B)}\,\Rt\,p\big(\B/\alpha_\B\big)\,\Rt^T,
\]
with $p\approx\delta/t$ as above (Theorem~\ref{thm:B} justifies the first equality).
The degree drops from about $\kappa_\alpha(\A)\log(1/\varepsilon)$ to about $\kappa_\alpha(\B)\log(1/\varepsilon)$.
Both are effective condition numbers, because the cost is always set by
the encoded matrix. For $\A$ we only use $\kappa_\alpha(\A)\ge\kappa(\A)$, with equality for an
ideal encoding; this is the best case for the unpreconditioned solver. For
$\B$ the gap between $\kappa_\alpha(\B)$ and $\kappa(\B)$ is exactly what the encoding lever
below controls, so $\kappa(\B)$ alone would be misleading. For
the Poisson control problem of Section~\ref{sec:poisson} with $d=1$ and
$h=2^{-8}$, this is $\kappa_\alpha(\A)\ge\kappa(\A)\approx2.6\cdot10^4$ against $\kappa_\alpha(\B)\approx5$--$10$. The
latter assumes an ideal encoding of the factor $G$ of
Section~\ref{sec:poisson}. There $\B=\G^T\J\G$ with $\G=I_2\otimes G$, and $G$ maps
multilevel coefficients to values and gradients at quadrature points. The
middle factor $\J$ costs nothing, so $\alpha_\B=\alpha_G^2$, which is $\|G\|^2$ in the ideal
case. The value for the explicit construction is still to be measured
(Section~\ref{sec:poisson}).

\emph{Two levers, one classical and one quantum.} The ratio
$\delta=\sigma_{\min}(\B)/\alpha_\B$ can be spoiled from either side.
\begin{enumerate}
\item \emph{Spectral (classical).} $\sigma_{\min}(\B)$ must be comparable to $\|\B\|$.
  This is ordinary spectral equivalence of $P$ and $\A$, i.e. eigenvalue
  inclusion intervals. In Section~\ref{sec:poisson} it is provided by the
  Sch\"oberl--Zulehner block structure combined with BPX.
\item \emph{Encoding (quantum).} $\alpha_\B$ must be comparable to $\|\B\|$. This
  fails if $\B$ is encoded as the product $\Rt^T\cdot\A\cdot\Rt$: then
  $\alpha_\B\gtrsim\|R\|^2\|\A\|$, which is of the order of $\kappa(\A)$ (Section~\ref{sec:poisson},
  \eqref{eq:naivealpha}), and lever~1 is undone. It succeeds if $\B$ is
  written as a product of factors that are each encoded close to their own
  norm, and whose norms multiply to about $\|\B\|$ (no cancellation between
  the factors). A Gram factorization $F^TF$ has the second property
  automatically, since $\|F^TF\|=\|F\|^2$. Our $\B=\G^T\J\G$ (Section~\ref{sec:poisson}) is
  of this type, with a reflection $\J$ in the middle to carry the
  indefiniteness.
\end{enumerate}

\emph{Comparison with classical preconditioning.}
\begin{itemize}
\item \emph{Where the preconditioner sits.} Classically, $P$ is applied once per
  iteration, usually by an inner solve such as a multigrid cycle. Here the
  preconditioner is built into the matrix that is encoded, so every call of
  the circuit already multiplies by the preconditioned operator. There are
  no inner solves: $P^{-1}$ must be available as a multiplication by a factor
  $\Rt$ (principle~(D3) below).
\item \emph{What is optimized.} A Krylov method chooses its polynomial
  adaptively, from the actual spectrum and the right-hand side, so
  clustering of eigenvalues helps. Here the polynomial is fixed in advance
  by the interval $[\delta,1]$, so only the interval counts (principle~(D1)).
\item \emph{$\kappa$ versus $\sqrt\kappa$.} For SPD systems, Chebyshev iteration or CG
  need about $\sqrt\kappa$ steps, because they minimize a residual polynomial
  normalized by $p(0)=1$, a relative error measure. QSVT must approximate
  the inverse itself uniformly on $[\delta,1]$, while staying bounded by one on
  $[-1,1]$. This needs degree of order $\kappa_\alpha$. The quantum solver therefore
  depends more strongly on the condition number than CG does, which makes
  preconditioning \emph{more} important, not less.
\item \emph{The ends of the product.} In practice $\Rt$ and $\Rt^T$ are not
  applied separately. The right-hand side is prepared as $\Rt^Tf=\G^Tg$, and
  outputs are read through $\G$ (principle~(D5)).
\end{itemize}

\paragraph{Design principles.} Equation~\eqref{eq:cost} changes what a good
preconditioner is.
\begin{enumerate}
\item[(D1)] \emph{Only $\kappa_\alpha$ counts.} QSVT does not adapt to the spectrum,
  unlike MINRES. Three distinct eigenvalues help MINRES but not a QLSA.
  What we need are inclusion intervals $[-b,-a]\cup[a,b]$ with $b/a$
  independent of the parameters, i.e. exactly the robustness results of
  \cite{SZ07,PW12}.
\item[(D2)] \emph{Subnormalizations count.} The preconditioned matrix must be
  encoded with $\alpha$ close to its norm. Multiplying separately encoded
  factors such as $P^{-1/2}$, $\A$ and $P^{-1/2}$ fails. This is why
  \cite{DP26} precondition a \emph{factor} of the stiffness matrix rather than
  the matrix itself; see also \cite{LS25}.
\item[(D3)] \emph{Preconditioners multiply.} A classical inner solve
  (e.g. a multigrid cycle) has to become a circuit that multiplies by
  $P^{-1}$ or by a factor of it. Additive multilevel methods \cite{BPX90,Xu92}
  are sums of products of structured sparse matrices and fit this
  model. Multiplicative multigrid does not.
\item[(D4)] \emph{Kernels are harmless.} QSVT inverts on the range, so
  redundant generating systems, such as multilevel ones, are allowed. This
  is used in \cite{DP26} for the BPX generating system.
\item[(D5)] \emph{Organize for QoIs.} Right-hand sides and outputs should be
  accessible without applying further badly scaled operators.
\end{enumerate}

\section{Distributed control of the Poisson equation}\label{sec:poisson}

\subsection{Problem and discretization}

Let $\Omega=(0,1)^d$, $d\le3$, with desired state $\hat y\in L^2(\Omega)$ and
regularization parameter $\beta>0$. We solve
\begin{equation}\label{eq:ocp}
\min_{y,u}\ \tfrac12\|y-\hat y\|_{L^2(\Omega)}^2+\tfrac{\beta}{2}\|u\|_{L^2(\Omega)}^2
\quad\text{s.t.}\quad -\Delta y=u\ \text{in }\Omega,\quad y=0\ \text{on }\partial\Omega .
\end{equation}

\paragraph{Grid and matrices.} We use $Q_1$ elements on a uniform grid with
$h=2^{-L}$. There are $2^L$ elements and $2^L-1$ interior nodes per direction,
so $n=(2^L-1)^d$ unknowns. In 1D,
$M_1=\frac h6\operatorname{tridiag}(1,4,1)$ and $K_1=\frac1h\operatorname{tridiag}(-1,2,-1)$. In $d$
dimensions the matrices are Kronecker products, e.g.
$M=M_1\otimes M_1$ and $K=K_1\otimes M_1+M_1\otimes K_1$ for $d=2$. This structure matches the
tensor-product structure of a qubit register.

\paragraph{Factorizations through quadrature points.} Let $\xi_{1,2}=\frac12\mp\frac1{2\sqrt3}$
be the two Gauss points on $[0,1]$. Let $F_0,F_1\in\R^{2^{L+1}\times(2^L-1)}$ map nodal
values to $\sqrt{h/2}$ times the values, respectively derivatives, of the finite
element function at the Gauss points of all elements. Two-point Gauss
quadrature is exact for the integrands involved, so
\begin{equation}\label{eq:fact1d}
M_1=F_0^TF_0,\qquad K_1=F_1^TF_1 .
\end{equation}
In $d$ dimensions let $F_{\mathrm{val}}=F_0^{\otimes d}$ and let $F_\nabla$ stack the factors
$F^{(i)}=F_0\otimes\cdots\otimes F_1\otimes\cdots\otimes F_0$ ($F_1$ in position $i$). Then
$M=F_{\mathrm{val}}^TF_{\mathrm{val}}$ and $K=F_\nabla^TF_\nabla$. The factors have
$2^{d(L+1)}$ rows, $2^d$ nonzeros per row and at most $4^d$ per column.

This is the quadrature-point version of the factorization of Deiml and
Peterseim \cite[(6.1)]{DP26}, which uses an $L^2$-orthonormal (Legendre) basis
of discontinuous $Q_1$ functions on each element instead of weighted values
at Gauss points. For $Q_1$ the two are related by an orthogonal $2\times2$ map per
element and direction, so the statements are equivalent. The quadrature
form includes the mass matrix, as in \cite{YYZ26}, and extends directly to
$Q_k$ and $P_k$ elements (Remark~\ref{rem:elements}).

\begin{remark}[Other elements]\label{rem:elements}
The construction is not specific to $Q_1$. Consider a nested family of
conforming Lagrange spaces: $Q_k$ on Cartesian grids, or $P_k$ on uniformly
red-refined, possibly unstructured simplicial meshes. Take a quadrature rule
with positive weights that is exact for squares of the local polynomial
space.
\begin{itemize}
\item The factorizations $M=F_{\mathrm{val}}^TF_{\mathrm{val}}$ and $K=F_\nabla^TF_\nabla$ are
  exact.
\item The quadrature points are unisolvent, since $0<\int p^2=\sum_qw_qp(x_q)^2$.
\item Lemma~\ref{lem:commute} below holds with an element-local isometry,
  the same for values and gradients, because derivatives stay in the local
  space.
\item For $Q_k$ the Kronecker form $M=M_1\otimes M_1$, $K=K_1\otimes M_1+M_1\otimes K_1$
  ($d=2$) holds for every $k$. The nodal basis functions are products
  $\phi_i(x)\phi_j(y)$ of 1D degree-$k$ basis functions, so the integrals
  factor; only $M_1$ and $K_1$ change. For $Q_2$ they are 2-periodic (vertex
  and midpoint nodes alternate), with $2^{L+1}-1$ unknowns per direction.
  Serendipity elements and $P_k$ on simplices are not tensor products.
\item Row counts that are not powers of two are zero-padded.
\item Variable coefficients enter as one diagonal factor.
\end{itemize}
Mass lumping, curved elements and non-nested hierarchies keep the
factorization only up to spectral equivalence. Unstructured meshes lose the
Kronecker structure of the factors, which then need sparse-access oracles.
We verified these statements numerically for $Q_2$ and $Q_3$ (1D and 2D) and
for $P_1$ and $P_2$ on triangles. All factorization, commuting and isometry
errors are below $3.3\cdot10^{-14}$.
\end{remark}

\paragraph{Discrete problem and KKT system.} Let $b=F_{\mathrm{val}}^T\hat y_q$, where
$\hat y_q$ holds the weighted values of $\hat y$ at the Gauss points. We minimize
$\frac12y^TMy-b^Ty+\frac\beta2u^TMu$ subject to $Ky=Mu$. The optimality
conditions are
\begin{equation}\label{eq:kkt3}
\begin{bmatrix} M & 0 & K\\ 0 & \beta M & -M\\ K & -M & 0\end{bmatrix}
\begin{bmatrix} y\\ u\\ p\end{bmatrix}=\begin{bmatrix} b\\ 0\\ 0\end{bmatrix}.
\end{equation}
We eliminate $u=p/\beta$, set $q=p/\sqrt\beta$ and scale the second equation.
This gives
\begin{equation}\label{eq:kkt2}
\A x=f,\qquad
\A=\begin{bmatrix} M & \sqrt\beta K\\ \sqrt\beta K & -M\end{bmatrix},\qquad
x=\begin{bmatrix} y\\ q\end{bmatrix},\quad f=\begin{bmatrix} b\\ 0\end{bmatrix},
\end{equation}
with $u=q/\sqrt\beta$. The matrix $\A$ is symmetric and indefinite, with
$\kappa(\A)=O(1+\sqrt\beta h^{-2})$ (Table~\ref{tab:poisson}). Classically, MINRES
with a robust block preconditioner and multigrid solves~\eqref{eq:kkt2} in
$O(n)$ operations \cite{RDW10,SZ07,PW12}.

\subsection{A block preconditioner with a single elliptic operator}

Let $A_\beta:=M+\sqrt\beta K$, the stiffness matrix of $I-\sqrt\beta\Delta$.

\begin{theorem}[\cite{SZ07}]\label{thm:sz}
For $P=\operatorname{blkdiag}(A_\beta,A_\beta)$,
$\operatorname{spec}(P^{-1}\A)\subset[-1,-1/\sqrt2]\cup[1/\sqrt2,1]$ for all $h$ and $\beta$.
\end{theorem}
\begin{proof}
With $X^TMX=I$ and $X^TKX=\operatorname{diag}(\lambda_j)$, the pencil $(\A,P)$ decouples into
$2\times2$ pencils $\big(\begin{bmatrix}1&\mu_j\\ \mu_j&-1\end{bmatrix},(1+\mu_j)I\big)$ with
$\mu_j=\sqrt\beta\lambda_j>0$. Their eigenvalues are
$\pm\sqrt{1+\mu_j^2}/(1+\mu_j)\in\pm[1/\sqrt2,1]$.
\end{proof}

\begin{remark}
In the Pearson--Wathen approximation of the Schur complement of
\eqref{eq:kkt3} \cite{PW12}, $\widehat S=(K+M/\sqrt\beta)M^{-1}(K+M/\sqrt\beta)$,
the operator to invert is again $K+M/\sqrt\beta=A_\beta/\sqrt\beta$. We use the
form of Theorem~\ref{thm:sz} because it avoids $M^{-1}$ and has the sharper
interval.
\end{remark}

\subsection{A multilevel factor and level weights}\label{sec:weights}

For levels $\ell=1,\dots,L$ with $h_\ell=2^{-\ell}$, let $P_\ell\in\R^{n\times n_\ell}$ be the
prolongation from level $\ell$ to level $L$ ($P_L=I$). The level matrices
$A_{\beta,\ell}$ have constant diagonal on a uniform grid:
\[
\operatorname{diag}(A_{\beta,\ell})=\delta_\ell I,\qquad
\delta_\ell=\big(\tfrac{2h_\ell}{3}\big)^d+\sqrt\beta\,\tfrac{2d}{h_\ell}\big(\tfrac{2h_\ell}{3}\big)^{d-1}.
\]
Following \cite[Sec.~5]{DP26}, we use BPX in factored form, as a generating
system $R$ with $C=RR^T$. We define the weighted, diagonally scaled BPX factor
\begin{equation}\label{eq:bpx}
R=\big[\,\omega_1P_1,\ \dots,\ \omega_LP_L\,\big],\qquad
C=RR^T=\sum_{\ell=1}^L\omega_\ell^2P_\ell P_\ell^T,\qquad
\omega_\ell^2=\frac{\theta_\ell}{\delta_\ell},
\end{equation}
with level weights
\begin{equation}\label{eq:theta}
\theta_\ell=\left(\frac{\min\big(1,\sqrt\beta\,h_\ell^{-2}\big)}{\min\big(1,\sqrt\beta\,h_L^{-2}\big)}\right)^{s},\qquad s\ge0 .
\end{equation}
The choice $s=0$ is plain BPX. Since $P_L=I$, $R$ has full row rank. The
denominator normalizes the finest level to $\theta_L=1$, so that $\theta_\ell\le1$
for all $\ell$. It changes $C$ only by a positive scalar and therefore does not
affect $\kappa(CA_\beta)$. It does matter for the interval $[c_0,c_1]$ below:
without it, $c_0$ would shrink like $(\sqrt\beta h_L^{-2})^s$ whenever the finest
grid does not resolve $\beta$. Since the cost bounds below involve $c_0$
itself, not only $c_1/c_0$, the normalization is needed.

\begin{assumption}\label{ass:bpx}
$\operatorname{spec}(CA_\beta)\subset[c_0,c_1]$, with $c_1/c_0$ bounded independently of $h$
and, ideally, of $\beta$.
\end{assumption}

For fixed $\beta$ and $s=0$, this is the classical BPX theory for
$H^1$-elliptic problems \cite{BPX90,Xu92,Oswald94}. Robustness in $\beta$
concerns the singularly perturbed operator $I-\sqrt\beta\Delta$.
\begin{itemize}
\item \emph{Plain BPX} is robust in $h$ but degrades as $\beta\to0$
  (Section~\ref{sec:num}): coarse levels on which the mass term dominates
  ($\sqrt\beta h_\ell^{-2}<1$) all count the same smooth modes.
\item \emph{The weights}~\eqref{eq:theta} with $s>0$ damp these levels
  geometrically. With $s=1$ they keep $\kappa(CA_\beta)\le6.2$ in all our tests in
  which the grid resolves $\beta$ ($h\lesssim\beta^{1/4}$). If it does not, then
  $A_\beta\approx M$, and $\kappa(CA_\beta)$ approaches the ``mass-matrix floor''
  $3^d/(1-4^{-s})$, i.e.\ $4$, $12$ and $36$ for $d=1,2,3$ and $s=1$.
\item \emph{Truncation}, i.e.\ removing the mass-dominated levels altogether
  (the limit $s\to\infty$), is also robust, but its constants are worse at
  intermediate $\beta$: for example $\kappa(CA_\beta)=19.2$ versus $6.1$ for $s=1$
  at $d=1$, $L=10$, $\beta=10^{-4}$.
\end{itemize}

\begin{theorem}[robustness of the weighted factor]\label{thm:robust}
Let $s>0$, $Q_1$ elements on uniform dyadic grids of $(0,1)^d$, and
homogeneous Dirichlet conditions. Then Assumption~\ref{ass:bpx} holds with
$c_0,c_1$ depending only on $d$ and $s$, uniformly in $L$ (i.e.\ $h$) and
$\beta>0$.
\end{theorem}
A self-contained proof with explicit constants is given in the
supplementary note~\cite{SuppProof}. It reduces $x^TC^{-1}x$ to a minimal
multilevel decomposition (the fictitious-space identity
\cite{GO95b,XZ02}), compares nodal and $L^2$ norms by Jacobi scaling, and
uses two classical ingredients: the Jackson-type multilevel norm
equivalence \cite[proof of Thm.~1]{GO95} and the strengthened
Cauchy--Schwarz inequality \cite[Lemma~6.1]{Xu92}. The level weights make
the mass contributions of the levels a geometric series, which removes
the factor $\min(L,\log(1/\beta))$ of plain BPX. That factor is genuine
for $s=0$ \cite[Rem.~4.5]{KXZ16}. The explicit constants are pessimistic by
two to three orders of magnitude. The proof explains the structure, not
the size of $\kappa(CA_\beta)$.

\paragraph{Relation to robust multilevel methods.} Additive multilevel
preconditioners for $I-\sqrt\beta\,\Delta$ that are uniform in $\beta$ and $h$ are
classical.
\begin{itemize}
\item Griebel and Oswald \cite[Thms.~1--2]{GO95} prove robustness for
  $Q_1$ elements on $(0,1)^d$, including nodal BPX and multilevel diagonal
  scaling, with the mass-dominated coarse levels dropped. This is the
  truncated limit $s\to\infty$ of~\eqref{eq:theta}.
\item Bramble, Pasciak and Vassilevski \cite{BPV00} use the level weights
  $(\sqrt\beta\,h_\ell^{-2}+1)^{-1}$ on differences of $L^2$ projections. Bacuta, Hayes and
  Jacavage \cite{BHJ21} give a nodal version with lumped mass. Rewritten
  in terms of the $P_\ell P_\ell^T$, these weights are within level-independent
  factors of ours with $s=1$ when the grid resolves $\beta$.
\item Yang, Yu and Zhang \cite[Thm.~2.1]{YYZ26} use our $s=0$ with exact
  inverse mass matrices. Their bound $12dL^2$ is uniform in the reaction
  coefficient but not in $L$.
\end{itemize}
Our contribution is therefore not robustness as such. It is the form
required by the block encoding: a Jacobi-scaled factor $R$ in which every
level is a scaled prolongation; a family of weights $s>0$ that contains
the known choices; the normalization $\theta_L=1$, which makes $[c_0,c_1]$
itself independent of $\beta$, which is needed in Section~\ref{sec:heat};
and explicit constants. A better mass factor than Jacobi scaling would
lower the floor $3^d/(1-4^{-s})$.

\subsection{The two-sided system and its factorization}

Let $\Rt=I_2\otimes R$ and
\begin{equation}\label{eq:B}
\B:=\Rt^T\A\Rt=\begin{bmatrix}R^TMR & \sqrt\beta R^TKR\\ \sqrt\beta R^TKR & -R^TMR\end{bmatrix}.
\end{equation}

\begin{theorem}\label{thm:B}
Under Assumption~\ref{ass:bpx}:
\begin{enumerate}
\item $\A x=f$ if and only if $x=\Rt z$ for a solution $z$ of the consistent
  system $\B z=\Rt^Tf$. In particular $x=\Rt\B^+\Rt^Tf$.
\item The nonzero eigenvalues of $\B$ lie in $[-c_1,-c_0/\sqrt2]\cup[c_0/\sqrt2,c_1]$.
\end{enumerate}
\end{theorem}
\begin{proof}
(1) $\Rt$ is surjective, so $\Rt^T$ is injective, and
$\operatorname{range}(\B)=\Rt^T\A\R^{2n}=\operatorname{range}(\Rt^T)$.
(2) The nonzero eigenvalues of $\Rt^T\A\Rt$ are those of
$\widetilde C^{1/2}\A\widetilde C^{1/2}=S(P^{-1/2}\A P^{-1/2})S^T$, where
$\widetilde C=I_2\otimes C$ and $S=\widetilde C^{1/2}P^{1/2}$. By Ostrowski's theorem
\cite[Thm.~4.5.9]{HJ13}, each one equals $\theta\lambda$, where $\lambda$ is an
eigenvalue as in Theorem~\ref{thm:sz} and $\theta\in\operatorname{spec}(SS^T)=\operatorname{spec}(I_2\otimes CA_\beta)
\subset[c_0,c_1]$.
\end{proof}

\paragraph{Factorization with a reflection.}

\emph{Why a factorization is needed.} Classically, a preconditioned operator
$\Rt^T\A\Rt$ is applied factor by factor, and the norms of the factors are
irrelevant. A Krylov method sees only the spectrum of the product. On a
quantum computer, a product of block encodings encodes
$\Rt^T\A\Rt/(\alpha_R^2\alpha_\A)$, because every factor is applied with a loss of
amplitude. Its subnormalization satisfies
\begin{equation}\label{eq:naivealpha}
\alpha\ \ge\ \|R\|^2\,\|\A\|\ =\ \|C\|\,\|\A\|\ \ge\ c_0\|A_\beta^{-1}\|\cdot\tfrac12\|A_\beta\|\ =\ \tfrac{c_0}{2}\,\kappa(A_\beta).
\end{equation}
The steps are as follows.
\begin{itemize}
\item Each factor $X$ has $\alpha_X\ge\|X\|$, since $X/\alpha_X$ is a block of a
  unitary matrix, and the $\alpha$'s multiply.
\item $\|R\|^2=\|RR^T\|=\|C\|\ge c_0\|A_\beta^{-1}\|$ by Assumption~\ref{ass:bpx}.
\item $\|A_\beta\|\le\|M\|+\sqrt\beta\|K\|\le2\|\A\|$, because $M$ and $\sqrt\beta K$ are
  blocks of $\A$.
\end{itemize}
Since $\sigma_{\min}(\B)\le\|\B\|\le c_1$, the effective condition number of the
product encoding is at least $\frac{c_0}{2c_1}\kappa(A_\beta)$, with
$\kappa(A_\beta)\sim1+\sqrt\beta h^{-2}$. Numerically this is close to $\kappa(\A)$. For
example, for $d=1$, $L=8$, $\beta=1$: $\|R\|^2\|\A\|=3.0\cdot10^4$, $\kappa(\A)=2.6\cdot10^4$, while
the factored form below gives $\|G\|^2/\sigma_{\min}(\B)=5.0$. QSVT then has to
invert on $[\sigma_{\min}/\alpha,1]$, with degree $\sim\alpha/\sigma_{\min}$. It is non-adaptive:
it behaves like Chebyshev semi-iteration with the upper eigenvalue bound
overestimated as $\alpha$ instead of $\lambda_{\max}$. So the cost would be of the
order of $\kappa(\A)$, as without preconditioning. Krylov methods do not
suffer from this, because they adapt to the actual spectrum.

\emph{What a good factorization needs.} We look for a factorization
$\B=F_1F_2\cdots F_m$ with two properties.
\begin{enumerate}
\item[(W1)] \emph{Tight factors:} each factor has an efficient block encoding
  with $\alpha(F_i)\approx\|F_i\|$.
\item[(W2)] \emph{No cancellation:} $\prod_i\|F_i\|\approx\|\B\|$. That is, the vectors on which
  one factor is large are not damped by its neighbours.
\end{enumerate}
Then $\alpha_\B=\prod_i\alpha(F_i)\approx\|\B\|$. The naive product $\Rt^T\cdot\A\cdot\Rt$ satisfies (W1)
but violates (W2): $R$ is large on smooth vectors, where $\A$ is small, and
$\|R\|^2\|\A\|$ exceeds $\|\B\|$ by a factor of order $\kappa(A_\beta)$.
\begin{itemize}
\item \emph{Gram factorizations satisfy (W2) automatically,} because
  $\|F^TF\|=\|F\|^2$ exactly. This is why we factor through
  $G=[F_{\mathrm{val}}R;\,\beta^{1/4}F_\nabla R]$, which has $G^TG=R^TA_\beta R$.
\item \emph{An indefinite $\B$ needs a reflection in the middle} (next
  paragraph). (W2) then holds up to a bounded factor: $\|\G\|^2\|\J\|=\|G\|^2\le c_1$,
  while $\|\B\|\ge c_0/\sqrt2$ by Theorem~\ref{thm:B}, so
  $\prod_i\|F_i\|/\|\B\|\le\sqrt2\,c_1/c_0$.
\item \emph{(W1)} is the subject of Section~\ref{sec:blockG}
  (Lemma~\ref{lem:commute} and Proposition~\ref{prop:alphaG}).
\end{itemize}

\emph{Why a reflection in the middle.} $\B$ is indefinite, so it cannot be
written as $G^TG$. Let
\[
G=\begin{bmatrix}F_{\mathrm{val}}R\\ \beta^{1/4}F_\nabla R\end{bmatrix},\qquad
\G=I_2\otimes G=\begin{bmatrix}G&0\\0&G\end{bmatrix},\qquad
\J=\begin{bmatrix}E_{\mathrm{val}} & E_\nabla\\ E_\nabla & -E_{\mathrm{val}}\end{bmatrix},
\]
where $E_{\mathrm{val}}$ and $E_\nabla=I-E_{\mathrm{val}}$ are the diagonal 0/1 matrices that
select the value rows and the gradient rows of $G$. Then
\begin{equation}\label{eq:BGJG}
\B=\G^T\J\G,\qquad \J=\J^T=\J^{-1},\qquad \|G\|^2=\|R^TA_\beta R\|\le c_1 .
\end{equation}
This is a $J$-factorization, as in $LDL^T$ with $D=\pm1$ or hyperbolic QR.
$G$ maps coefficients to values $v$ and gradients $g$ at the quadrature
points, for both $y$ and $q$. $\J$ is the matrix of the bilinear form
\[
z^T\B z=\|v_y\|^2-\|v_q\|^2+2\langle g_y,g_q\rangle .
\]
On value rows it multiplies the $y$-block by $+1$ and the $q$-block by
$-1$, giving $\pm R^TMR$. On gradient rows it swaps the two blocks, giving
the off-diagonal $\sqrt\beta R^TKR$.
\begin{itemize}
\item $\J$ is a signed permutation matrix. On the device it is a single
  two-qubit gate acting on the qubit that labels the block ($y$ or $q$)
  and the qubit that labels the row type (value or gradient): ``if value
  row, flip the sign of the $q$-block; if gradient row, swap the
  blocks''. Its subnormalization is $1$.
\item The right-hand side is $\Rt^Tf=\G^Tg$, where $g$ is $\hat y_q$ placed in
  the value rows of the $y$-block (zero elsewhere).
\item Outputs are accessed through $\G$ as well, since $G_{\mathrm{val}}z_y=F_{\mathrm{val}}y$.
  For example, $\langle w,y\rangle_M$ is an overlap of two states prepared
  with $\G$. So $R$ is never applied on its own (D5).
\end{itemize}

\subsection{Block encoding of the multilevel factor}\label{sec:blockG}

A product encoding of $F_\nabla P_\ell$ loses a factor $h_\ell/h$. The reason is
that $\|F_\nabla\|$ is attained by oscillatory vectors and $\|P_\ell\|$ by smooth
ones. Numerically, the loss reaches $148$ for $L=8$ (Table~\ref{tab:loss}).
The remedy is to differentiate on the coarse level and only then transfer
to the fine level.

\begin{lemma}[commuting with prolongation; cf.\ {\cite[Thm.~6.2]{DP26}}]\label{lem:commute}
For each level $\ell$ there is an isometry $T_\ell$ ($T_\ell^TT_\ell=I$) with
\[
F_{\mathrm{val}}P_\ell=T_\ell F_{\mathrm{val},\ell},\qquad F^{(i)}P_\ell=T_\ell F^{(i)}_\ell\quad(i=1,\dots,d).
\]
$T_\ell$ is element-local and the same on every coarse element. In 1D,
$T_\ell=I_{2^\ell}\otimes\widehat T_{L-\ell}$, where $\widehat T_k\in\R^{2^{k+1}\times2}$ maps the two
weighted Gauss values of an affine function on a coarse element to its
weighted values at the $2^{k+1}$ fine Gauss points. In $d$ dimensions,
$T_\ell$ is the $d$-fold Kronecker product.
\end{lemma}
\begin{proof}
On a coarse element, coarse functions and their partial derivatives are
affine in each direction. The two Gauss points are unisolvent for affine
functions, which defines $\widehat T_k$. Both the coarse and the fine weighted
values compute the exact $L^2$ norm of the same affine function, because
two-point Gauss quadrature is exact for quadratics. Hence $\widehat T_k$ is an
isometry.
\end{proof}

For $Q_1$ and the gradient part, Lemma~\ref{lem:commute} is the commutative
diagram of Deiml and Peterseim \cite[Thm.~6.2]{DP26}: differentiate on the
coarse level, then include the coarse discontinuous space into the fine one
by a map with orthonormal columns. The finite-difference analogue is
\cite[Lemma~4.3]{JLMY25}. The version with values and gradients, i.e. with
the mass term, is \cite[Lemma~3.3]{YYZ26}. Our contribution here is the
quadrature formulation, which covers $Q_k$ and $P_k$ elements
(Remark~\ref{rem:elements}).

Some matrix $T$ with $F P_\ell=TF_\ell$ always exists when $F_\ell$ has full column
rank, e.g. $T=FP_\ell F_\ell^+$. The content of Lemma~\ref{lem:commute} is that
$T_\ell$ is an isometry and element-local. This is what gives a cheap block
encoding with subnormalization one. By the lemma,
\begin{equation}\label{eq:Gsum}
G=\big[\,T_1G_1,\ \dots,\ T_LG_L\,\big],\qquad
G_\ell=\omega_\ell\begin{bmatrix}F_{\mathrm{val},\ell}\\ \beta^{1/4}F^{(1)}_\ell\\ \vdots\\ \beta^{1/4}F^{(d)}_\ell\end{bmatrix},\qquad
G_\ell^TG_\ell=\theta_\ell\delta_\ell^{-1}A_{\beta,\ell} .
\end{equation}

\begin{proposition}\label{prop:alphaG}
$G$ has a block encoding with subnormalization $\alpha_G\le C_d\sqrt L$, where
$C_d$ depends only on $d$.
\end{proposition}
\begin{proof}[Sketch]
\begin{itemize}
\item $G_\ell$ has $2^d$ nonzeros per row and $(d+1)4^d$ per column.
\item By the definition of $\delta_\ell$, its entries are bounded by $1$,
  uniformly in $h$ and $\beta$. The sparse-access encoding therefore gives
  $\alpha(G_\ell)\le(2^d(d+1)4^d)^{1/2}$.
\item Each $T_\ell$ is an isometry, so it can be completed to a unitary.
\item Concatenating the $L$ levels costs a factor $\sqrt L$.
\end{itemize}
\end{proof}
Since $\B=\G^T\J\G$, the subnormalization of $\B$ is $\alpha_G^2=O(L)$. This matches
the linear growth in $L$ that Deiml and Peterseim obtain for the
preconditioned stiffness factor \cite[Thm.~6.3]{DP26}.
The sparse-access constant is pessimistic ($C_3\approx45$). Encoding $G_\ell$ as
a weighted stack of Kronecker products of 1D two-term factors should give
constants close to $\|G_\ell\|\le\sqrt{\theta_\ell}\,3^{d/2}$.

\paragraph{From bound to measurement.} The gap between the lower bound
$\|G\|\approx2$ (Table~\ref{tab:poisson}) and the upper bound $C_d\sqrt L$ matters,
because $\alpha_G$ enters the cost squared. For $d=3$ the two bounds differ by
a factor of about $2000$ in $\alpha_G^2$. We therefore plan to build the block
encoding of Proposition~\ref{prop:alphaG} as an explicit unitary matrix
for small grids ($d=1$ with $L\le6$, $d=2$ with $L\le4$), using exactly the
factors a circuit would use:
\begin{itemize}
\item LCUs of shifts for the 1D two-term factors, as in the example of
  Section~\ref{sec:prelim};
\item Kronecker products;
\item the weighted stack over values and gradients;
\item the isometries $T_\ell$, completed to unitaries;
\item the concatenation over levels.
\end{itemize}
For this matrix we check two things:
\begin{enumerate}
\item it is unitary, and its leading block equals $G/\alpha_G$ (correctness);
\item the value of $\alpha_G$ as a function of $L$, $d$ and $\beta$, compared with
  $\|G\|$ (cost). This shows whether the $\sqrt L$ appears and how large $C_d$
  really is.
\end{enumerate}
The effective condition number that QSVT faces is then
$\kappa_\alpha=\alpha_G^2/\sigma_{\min}(\B)$, not the ideal $\|G\|^2/\sigma_{\min}(\B)$.
\todo{carry out; can the $\sqrt L$ be avoided?}

\begin{theorem}\label{thm:main}
Under Assumption~\ref{ass:bpx}, QSVT applied to $\B=\G^T\J\G$ solves
$\B z=\G^Tg$ to accuracy $\varepsilon$, in the sense of~\eqref{eq:eps}, using
\[
O\Big(\frac{C_d^2L}{c_0}\log\frac1\varepsilon\Big)
\]
calls to block encodings of $G$, $G^T$ and $\J$. Each call costs
$\mathrm{poly}(L)$ gates. The count depends on $h$ only through $L=\log_2h^{-1}$, and
on $\beta$ only through $c_0$.
\end{theorem}
\begin{proof}
By~\eqref{eq:BGJG}, Theorem~\ref{thm:B} and Proposition~\ref{prop:alphaG},
$\kappa_\alpha(\B)\le\sqrt2\,\alpha_G^2/c_0\le\sqrt2\,C_d^2L/c_0$. The result follows
from~\eqref{eq:cost}.
\end{proof}
\noindent\todo{End-to-end statement: preparation of $g$, amplitude
amplification of the success probability, estimation of $\|z\|$, and readout
of the QoI.}

\section{The time-periodic heat equation}\label{sec:heat}

\subsection{Problem, discretization and Fourier transform in time}

We solve
\begin{equation}\label{eq:heatocp}
\min_{y,u}\ \frac12\int_0^T\!\!\|y-\hat y\|_{L^2(\Omega)}^2dt+\frac\beta2\int_0^T\!\!\|u\|_{L^2(\Omega)}^2dt
\quad\text{s.t.}\quad \partial_ty-\Delta y=u,\quad y|_{\partial\Omega}=0,\quad y(\cdot,0)=y(\cdot,T).
\end{equation}
Classically, this problem is treated with multiharmonic finite elements and
robust block preconditioners \cite{KK12,KKLWZ13}.
\begin{itemize}
\item \emph{Space:} $Q_1$ as in Section~\ref{sec:poisson}.
\item \emph{Time:} $N_t=2^{L_t}$ steps of length $\tau=T/N_t$ with periodic
  backward Euler, $M(y_k-y_{k-1})/\tau+Ky_k=Mu_k$, with $y_0=y_{N_t}$.
\end{itemize}
With the cyclic shift $\Sigma$ and $D_t=(I-\Sigma)/\tau$, the constraint reads
\[
\mathcal Ky=\mathbb Mu,\qquad \mathcal K=D_t\otimes M+I\otimes K,\qquad \mathbb M=I\otimes M .
\]
We approximate the objective by the rectangle rule. As
in~\eqref{eq:kkt2}, we eliminate $u$ and set $q=p/(\tau\sqrt\beta)$. This
gives the symmetric system
\begin{equation}\label{eq:heatkkt}
\A=\begin{bmatrix}\mathbb M & \sqrt\beta\,\mathcal K^T\\ \sqrt\beta\,\mathcal K & -\mathbb M\end{bmatrix},\qquad
x=\begin{bmatrix}y\\ q\end{bmatrix},\qquad f=\begin{bmatrix}\mathbb M\hat y\\ 0\end{bmatrix},\qquad
u=\frac{q}{\sqrt\beta}.
\end{equation}

$D_t$ is circulant: $D_t=\mathcal F^*\Lambda\mathcal F$, with the unitary DFT matrix $\mathcal F$
and
\begin{equation}\label{eq:lamk}
\lambda_k=\frac{1-e^{-2\pi ik/N_t}}{\tau}=\frac2\tau\sin\Big(\frac{\pi k}{N_t}\Big)e^{\,i(\pi/2-\pi k/N_t)},
\qquad\operatorname{Re}\lambda_k\ge0 .
\end{equation}
Applying $\mathcal F$ to the time index block-diagonalizes $\A$ into the Hermitian
blocks
\begin{equation}\label{eq:Ak}
\A_k=\begin{bmatrix}M & \sqrt\beta(K+\bar\lambda_kM)\\ \sqrt\beta(K+\lambda_kM) & -M\end{bmatrix},\qquad k=0,\dots,N_t-1 .
\end{equation}
On a quantum computer, $\mathcal F$ is the QFT on the $L_t$ time qubits, with
$O(L_t^2)$ gates. It is unitary, so it does not change the condition number
or the subnormalization. The analysis below uses only $\operatorname{Re}\lambda_k\ge0$. It
therefore also covers multiharmonic discretizations, where $\lambda_k=i\omega_k$.

\subsection{A robust preconditioner per frequency}

\begin{theorem}\label{thm:heatsz}
Let $r_k=1+\sqrt\beta|\lambda_k|$, $A_k=r_kM+\sqrt\beta K$ and $P_k=\operatorname{blkdiag}(A_k,A_k)$.
Then for all $h$, $\tau$, $\beta$ and $k$,
\[
\operatorname{spec}(P_k^{-1}\A_k)\subset[-1,-1/\sqrt3]\cup[1/\sqrt3,1].
\]
\end{theorem}
\begin{proof}
With $X^TMX=I$ and $X^TKX=\operatorname{diag}(\mu_j)$, the pencil decouples into
$2\times2$ pencils with eigenvalues $\pm\sqrt{1+x^2}/(1+c)$, where
$x=\sqrt\beta|\mu_j+\lambda_k|$ and $c=\sqrt\beta(\mu_j+|\lambda_k|)$.
\begin{itemize}
\item $c\ge x$, by the triangle inequality.
\item $c\le\sqrt2\,x$, because $\mu_j+|\lambda_k|\le(\mu_j+\operatorname{Re}\lambda_k)+|\operatorname{Im}\lambda_k|\le\sqrt2|\mu_j+\lambda_k|$
  when $\operatorname{Re}\lambda_k\ge0$.
\end{itemize}
Hence $\sqrt{1+x^2}/(1+\sqrt2x)\le\sqrt{1+x^2}/(1+c)\le1$. The lower bound is
minimal at $x=\sqrt2$, where it equals $1/\sqrt3$.
\end{proof}

\begin{remark}
\begin{itemize}
\item \emph{Time domain.} In the time domain the preconditioner is
  $\operatorname{blkdiag}(\mathcal P_t,\mathcal P_t)$ with
  $\mathcal P_t=I\otimes(M+\sqrt\beta K)+\sqrt\beta|D_t|\otimes M$ and $|D_t|=(D_t^TD_t)^{1/2}$.
  It is real and symmetric, although $\mathcal K$ is not. $|D_t|$ is the
  discrete analogue of $|\partial_t|$, which underlies the $H^{1/2}$-in-time
  norms of parabolic theory.
\item \emph{Relation to prior work.} Theorem~\ref{thm:heatsz} is the
  analogue, for this discretization, of the robust multiharmonic
  preconditioners of \cite{KK12,KKLWZ13}. \todo{compare constants and
  attribute.}
\end{itemize}
\end{remark}

\subsection{Each frequency is a Poisson problem}

We can write $A_k=r_k(M+\sqrt{\beta_k}K)$ with $\sqrt{\beta_k}=\sqrt\beta/r_k$. Each
frequency is therefore the operator of Section~\ref{sec:poisson} with
effective parameter $\beta_k$, scaled by $r_k$.
\begin{itemize}
\item \emph{Multilevel factor.} $R_k=[\omega_{1,k}P_1,\dots,\omega_{L,k}P_L]$ with
  $\omega_{\ell,k}^2=\theta_{\ell,k}/\delta_{\ell,k}$, where
  $\delta_{\ell,k}=r_k(\frac{2h_\ell}3)^d+\sqrt\beta\frac{2d}{h_\ell}(\frac{2h_\ell}3)^{d-1}$ and
  $\theta_{\ell,k}$ the normalized weights~\eqref{eq:theta} with $\beta_k$ in place
  of $\beta$. Then $C_kA_k=C(\beta_k)A_{\beta_k}$, where $C(\beta_k)$ and
  $A_{\beta_k}=M+\sqrt{\beta_k}K$ are the Poisson quantities for parameter $\beta_k$.
  By Theorem~\ref{thm:robust}, $\operatorname{spec}(C_kA_k)\subset[c_0,c_1]$ with the
  \emph{same} $c_0,c_1$ for all $k$ \cite{SuppProof}. This is where the
  normalization $\theta_{L,k}=1$ is essential: with unnormalized weights the
  interval would shrink by $\min(1,\sqrt{\beta_k}h_L^{-2})^s$, which is small for
  high frequencies when $\tau\lesssim2h_L^2$.
\item \emph{Frequency enters only through scalars.} The prolongations $P_\ell$
  and isometries $T_\ell$ are independent of $k$.
\item \emph{The level weights are no longer optional.} $\beta_k$ decreases from
  $\beta$ ($k=0$) to $\beta\tau^2/(\tau+2\sqrt\beta)^2$ ($k=N_t/2$), so high frequencies
  are mass-dominated.
\end{itemize}

\paragraph{Factorization with a scaled reflection.} Let
\[
G_k=\begin{bmatrix}\sqrt{r_k}\,F_{\mathrm{val}}R_k\\ \beta^{1/4}F_\nabla R_k\end{bmatrix},\qquad
G_k^*G_k=R_k^TA_kR_k .
\]
Then $\B_k:=\Rt_k^T\A_k\Rt_k=\G_k^*\J_k\G_k$, where
\begin{equation}\label{eq:Jk}
\J_k=\begin{bmatrix}j_{11}E_{\mathrm{val}} & j_{12}E_{\mathrm{val}}+E_\nabla\\ j_{21}E_{\mathrm{val}}+E_\nabla & j_{22}E_{\mathrm{val}}\end{bmatrix},\qquad
J_k=\begin{bmatrix}j_{11}&j_{12}\\ j_{21}&j_{22}\end{bmatrix}
=\frac1{r_k}\begin{bmatrix}1&\sqrt\beta\bar\lambda_k\\ \sqrt\beta\lambda_k&-1\end{bmatrix}=s_kU_k .
\end{equation}
On value rows, $\J_k$ applies the $2\times2$ matrix $J_k$ across the two blocks. On
gradient rows, it swaps the blocks as before. $U_k$ is a Hermitian unitary
$2\times2$ matrix (a reflection) with phase $\arg\lambda_k=\pi/2-\pi k/N_t$, and
\[
s_k=\frac{\sqrt{1+\beta|\lambda_k|^2}}{1+\sqrt\beta|\lambda_k|}\in\big[\tfrac1{\sqrt2},1\big].
\]
Hence $\|\J_k\|\le1$, and $\J_k$ is block encoded with subnormalization $1$.
The weight $\sqrt{r_k}$ on the value rows is essential: without it,
$\|J_k\|$ would grow like $\sqrt\beta/\tau$.

\begin{theorem}\label{thm:heatB}
Use the normalized weights, so that Assumption~\ref{ass:bpx} holds for all
$A_k$ with constants $c_0,c_1$ independent of $k$ (Theorem~\ref{thm:robust}). Then the nonzero eigenvalues of $\B=\operatorname{blkdiag}_k(\B_k)$ lie
in $[-c_1,-c_0/\sqrt3]\cup[c_0/\sqrt3,c_1]$, and
$\kappa_\alpha(\B)\le\sqrt3\,\alpha_G^2/c_0$ with $\alpha_G\le C_d\sqrt L$. In particular,
neither $N_t$ nor $\tau$ enters $\kappa_\alpha$.
\end{theorem}
\begin{proof}
Argue as for Theorem~\ref{thm:B} and Proposition~\ref{prop:alphaG},
frequency by frequency, with Theorem~\ref{thm:heatsz} in place of
Theorem~\ref{thm:sz}.
\end{proof}

QSVT applies $\B^+$ to all frequencies simultaneously, in superposition.
Compared with the Poisson problem, the cost per solve increases only by two
QFTs and the frequency-controlled scalars (Section~\ref{sec:impl}).

\section{Implementation}\label{sec:impl}

\paragraph{Registers.}
\begin{itemize}
\item $d(L+1)$ qubits for quadrature rows (element index and Gauss point
  per direction).
\item $dL$ qubits for nodal unknowns. Level-$\ell$ data use the leading $\ell$
  bits per direction. The unused index $2^\ell-1$, a boundary node, is mapped
  to zero by a comparator flag.
\item $\lceil\log_2L\rceil$ qubits for the level, $\lceil\log_2(d+1)\rceil$ for the
  component $c$, one qubit $b$ for the block, and $L_t$ qubits for time or
  frequency.
\item Ancillas for the block encodings and QSVT.
\end{itemize}
Coarse element indices are the leading bits of fine element indices, so
$T_\ell=I\otimes\widehat T_{L-\ell}$ acts on trailing bits only. Unused index
combinations lie in the kernel of $G$, which is harmless by (D4).

\paragraph{Circuits.}
\begin{enumerate}
\item \emph{1D two-term factors} $F_{0,\ell}$ and $F_{1,\ell}$: a shift by one on $\ell$
  qubits ($O(\ell)$ gates), an LCU of two diagonal patterns, and the boundary
  flag. Kronecker products and a weighted stack over $c$ then give $G_\ell$.
\item \emph{The isometry $\widehat T_k$}: its two orthonormal columns are the
  uniform vector (Hadamard gates) and a normalized discrete ramp with
  Gauss offsets. \todo{reference for linear-amplitude state preparation.}
\item \emph{Levels}: $G_\ell$ and $T_\ell$ controlled on the level register,
  followed by Hadamard gates and postselection on the level register.
\item \emph{$\J$}: one two-qubit gate (Poisson). For the heat equation,
  $\J_k$ needs $U_k$ controlled on $c$ and $k$, plus one ancilla rotation for
  $s_k$.
\item \emph{QSVT phase factors} for an odd polynomial approximating $1/t$
  on $[\sigma_{\min}/\alpha,1]$, of degree $O(\kappa_\alpha\log(1/\varepsilon))$. They are
  computed classically, e.g. as in \cite{DMWL21}, using a lower bound for
  $c_0$.
\item \emph{Frequency scalars (heat).} $\arg\lambda_k$ is affine in $k$, so it
  needs only single-qubit phase gates. The magnitude
  $|\lambda_k|=\frac2\tau\sin(\pi k/N_t)$ is an LCU of $\operatorname{diag}(e^{\pm i\pi k/N_t})$. The
  derived quantities $r_k^{\pm1/2}$, $\omega_{\ell,k}$ and $s_k$ come from QSVT
  polynomials or reversible arithmetic. \todo{compare costs.}
\item \emph{Right-hand side and readout.} Prepare $g$ (weighted values of
  $\hat y$ at Gauss points, and time steps for the heat problem); for smooth
  $\hat y$, efficient circuits exist \todo{refs}. Then apply $\G^T$, after a QFT
  in time for the heat problem. QoIs are overlaps $\langle\G g_w,\G z\rangle$,
  estimated by Hadamard tests and amplitude estimation; $\|z\|$ follows from
  the success probability.
\end{enumerate}

\paragraph{Classical precomputation.} All weights are in closed form. The
only nontrivial precomputation is the QSVT phase factors and the constant
$c_0$ that enters them.

\section{Numerical experiments}\label{sec:num}

All experiments use dense or sparse linear algebra in double precision,
with scripts in the folder \texttt{preconditioning/code} of the accompanying
repository:
\begin{itemize}
\item \texttt{qp\_poisson\_check.py}: Table~\ref{tab:poisson}, Table~\ref{tab:loss};
\item \texttt{qp\_bpx\_weights.py}: Table~\ref{tab:weights};
\item \texttt{qp\_bpx\_sparse.py}: Table~\ref{tab:3d}, extreme eigenvalues by
  Lanczos with CG-based inner solves;
\item \texttt{qp\_heat\_periodic\_check.py}: Table~\ref{tab:heat}.
\end{itemize}
Condition numbers on the range are ratios of the largest to the smallest
nonzero eigenvalue magnitudes.

\paragraph{Poisson control: preconditioner and factorization.}
Table~\ref{tab:poisson} shows:
\begin{itemize}
\item $\kappa(\A)$ grows like $\sqrt\beta h^{-2}$;
\item the interval of Theorem~\ref{thm:sz} is attained;
\item $\kappa(\B)$ on the range follows $\kappa(CA_\beta)$, as predicted by
  Theorem~\ref{thm:B};
\item $\|G\|$ is $O(1)$. This is the smallest possible $\alpha_G$; the value
  achieved by the explicit construction is still to be measured
  (Section~\ref{sec:poisson}, ``From bound to measurement'');
\item the factorization $\B=\G^T\J\G$ holds to machine precision in all
  cases.
\end{itemize}

\begin{table}[ht]
\centering\small
\begin{tabular}{rrrrcccc}
\toprule
$d$ & $L$ & $\beta$ & $\kappa(\A)$ & $|\operatorname{spec}(P^{-1}\A)|$ & $\kappa(CA_\beta)$ & $\kappa(\B)$ & $\|G\|$\\
\midrule
1 & 4 & $1$       & $1.0\cdot10^2$ & $[0.913,1.000]$ & 3.45 & 3.49 & 1.86\\
1 & 6 & $1$       & $1.7\cdot10^3$ & $[0.913,1.000]$ & 4.36 & 4.38 & 2.09\\
1 & 8 & $1$       & $2.6\cdot10^4$ & $[0.913,1.000]$ & 4.98 & 4.99 & 2.23\\
1 & 8 & $10^{-4}$ & $2.6\cdot10^3$ & $[0.708,1.000]$ & 4.43 & 4.50 & 2.11\\
1 & 8 & $10^{-8}$ & $2.6\cdot10^1$ & $[0.707,0.999]$ & 7.31 & 8.03 & 2.70\\
2 & 5 & $1$       & $2.1\cdot10^2$ & $[0.953,1.000]$ & 4.06 & 4.04 & 1.75\\
2 & 5 & $10^{-4}$ & $4.0\cdot10^1$ & $[0.707,0.996]$ & 3.96 & 3.47 & 1.74\\
2 & 5 & $10^{-6}$ & $4.2\cdot10^0$ & $[0.707,0.981]$ & 6.40 & 7.90 & 2.34\\
\bottomrule
\end{tabular}
\caption{Poisson control with plain BPX ($s=0$). $\kappa(\B)$ is computed on the
range. $\|G\|$ is a lower bound for the subnormalization $\alpha_G$ of any block
encoding of $G$.}
\label{tab:poisson}
\end{table}

\paragraph{Robustness in $h$ and $\beta$, including 3D.}
Tables~\ref{tab:weights} and~\ref{tab:3d} compare plain and weighted BPX.
\begin{itemize}
\item For fixed $\beta$, $\kappa(CA_\beta)$ is bounded in $h$ in all dimensions.
  It levels off slowly (3D, $\beta=1$: $4.2$, $5.3$, $6.1$ for $L=3,4,5$).
\item Plain BPX degrades as $\beta\to0$, up to $44.6$ in 3D.
\item The weights~\eqref{eq:theta} with $s=1$ keep $\kappa(CA_\beta)\le6.2$ whenever
  the grid resolves $\beta$ ($h\lesssim\beta^{1/4}$). In unresolved cases they
  approach the mass-matrix floor $3^d/(1-4^{-s})$ (4 in 1D and 12 in 2D for
  $s=1$). The worst value in 3D is $20.5$; for finer grids and tiny $\beta$ it
  should grow towards the 3D floor $36$.
\end{itemize}

\begin{table}[ht]
\centering\small
\begin{tabular}{rrrrrrrr}
\toprule
 & & \multicolumn{6}{c}{$\beta$}\\
\cmidrule(l){3-8}
$d,L$ & $s$ & $1$ & $10^{-2}$ & $10^{-4}$ & $10^{-6}$ & $10^{-8}$ & $10^{-10}$\\
\midrule
1, 10 & 0   & 5.42 & 5.32 & 5.02 & 4.92 & 7.33 & 9.80\\
1, 10 & 1/2 & 5.42 & 5.30 & 4.97 & 4.55 & 3.88 & 2.86\\
1, 10 & 1   & 5.42 & 5.33 & 6.14 & 5.37 & 4.92 & 3.61\\
\midrule
2, 4  & 0   & 3.58 & 3.49 & 3.62 & 7.23 & 18.89 & 26.31\\
2, 4  & 1/2 & 3.58 & 3.33 & 2.29 & 3.12 & 9.84 & 14.14\\
2, 4  & 1   & 3.58 & 3.29 & 2.26 & 1.97 & 7.16 & 10.51\\
\bottomrule
\end{tabular}
\caption{$\kappa(CA_\beta)$ for the level weights~\eqref{eq:theta}, including the
intermediate exponent $s=\tfrac12$. In 2D with $L=4$, the grid does not
resolve $\beta\le10^{-8}$.}
\label{tab:weights}
\end{table}

\begin{table}[ht]
\centering\small
\begin{tabular}{rrrcccccc}
\toprule
$d$ & $L$ & $n$ & $s$ & $\beta=1$ & $10^{-2}$ & $10^{-4}$ & $10^{-6}$ & $10^{-8}$\\
\midrule
2 & 6 & 3969  & 0 & 4.44 & 4.36 & 4.22 & 7.14 & 11.27\\
2 & 6 & 3969  & 1 & 4.44 & 4.27 & 3.53 & 2.83 & 1.64\\
\midrule
3 & 3 & 343   & 0 & 4.16 & 4.23 & 5.95 & 17.36 & 35.95\\
3 & 3 & 343   & 1 & 4.16 & 3.71 & 2.77 & 8.80 & 20.47\\
3 & 4 & 3375  & 0 & 5.30 & 5.31 & 6.41 & 15.75 & 44.55\\
3 & 4 & 3375  & 1 & 5.30 & 4.90 & 3.39 & 4.26 & 17.69\\
3 & 5 & 29791 & 0 & 6.08 & 6.05 & 6.65 & 14.19 & --\\
3 & 5 & 29791 & 1 & 6.08 & 5.74 & 4.46 & 3.12 & --\\
\bottomrule
\end{tabular}
\caption{$\kappa(CA_\beta)$ on larger grids, plain ($s=0$) and weighted ($s=1$)
BPX, computed by Lanczos.}
\label{tab:3d}
\end{table}

\paragraph{Subnormalization loss of the naive product.} Table~\ref{tab:loss}
gives the ratio $\|F_1\|\|P_\ell\|/\|F_1P_\ell\|$. A product encoding of $F_1P_\ell$
loses this factor. It equals $2^{L-\ell}$, up to a correction on the
coarsest levels. Lemma~\ref{lem:commute} removes it completely: the
commuted factors satisfy $\|T_\ell\|\,\|F_{1,\ell}\|=\|F_1P_\ell\|$.

\begin{table}[ht]
\centering\small
\begin{tabular}{lrrrrrrr}
\toprule
$\ell$ & 1 & 2 & 3 & 4 & 5 & 6 & 7\\
\midrule
$\|F_1\|\|P_\ell\|/\|F_1P_\ell\|$ & 147.8 & 65.8 & 32.2 & 16.0 & 8.0 & 4.0 & 2.0\\
\bottomrule
\end{tabular}
\caption{Loss of the naive product encoding, $d=1$, $L=8$.}
\label{tab:loss}
\end{table}

\paragraph{Time-periodic heat control.} Table~\ref{tab:heat} shows:
\begin{itemize}
\item the global eigenvalues of $\mathcal P^{-1}\A$ lie in the interval
  $[1/\sqrt3,1]$ of Theorem~\ref{thm:heatsz};
\item per frequency, the factorization $\B_k=\G_k^*\J_k\G_k$ holds and
  $\|J_k\|\le1$;
\item the per-frequency BPX constants behave as in the Poisson case with
  effective parameters $\beta_k$.
\end{itemize}

\begin{table}[ht]
\centering\small
\begin{tabular}{rrrrrc}
\toprule
$d$ & $L$ & $N_t$ & $\beta$ & $\kappa(\A)$ & $|\operatorname{spec}(\mathcal P^{-1}\A)|$\\
\midrule
1 & 5 & 32 & $1$       & $4.1\cdot10^2$ & $[0.714,1.000]$\\
1 & 5 & 32 & $10^{-2}$ & $2.9\cdot10^2$ & $[0.603,0.999]$\\
1 & 5 & 32 & $10^{-4}$ & $4.1\cdot10^1$ & $[0.666,0.992]$\\
1 & 5 & 32 & $10^{-6}$ & $4.2\cdot10^0$ & $[0.702,0.990]$\\
\midrule
2 & 3 & 16 & $1$       & $1.3\cdot10^1$ & $[0.811,0.999]$\\
2 & 3 & 16 & $10^{-2}$ & $1.2\cdot10^1$ & $[0.669,0.993]$\\
2 & 3 & 16 & $10^{-4}$ & $2.7\cdot10^0$ & $[0.684,0.936]$\\
2 & 3 & 16 & $10^{-6}$ & $4.4\cdot10^0$ & $[0.704,0.981]$\\
\midrule
\multicolumn{6}{l}{\todo{per-frequency BPX and factorization checks}}\\
\bottomrule
\end{tabular}
\caption{Time-periodic heat control. The bound of Theorem~\ref{thm:heatsz} is
$1/\sqrt3\approx0.577$.}
\label{tab:heat}
\end{table}

\section{Complexity and discussion}\label{sec:discussion}

\paragraph{Accuracies.} Let $\mathrm{tol}$ be the target accuracy of a QoI. It is
split into three parts:
\begin{itemize}
\item the discretization error;
\item the solver accuracy $\varepsilon$ of~\eqref{eq:eps};
\item the readout error $\eta$ of Section~\ref{sec:prelim}.
\end{itemize}
Each part is chosen of order $\mathrm{tol}$. Up to constants, the cost of a QoI to
accuracy $\mathrm{tol}$ is therefore (calls per solve with $\varepsilon\sim\mathrm{tol}$) times
(readout repetitions $\sim\mathrm{tol}^{-1}$).

\paragraph{Poisson control.} Suppose the discretization error of the QoI is
$O(h^2)$, so $h\sim\mathrm{tol}^{1/2}$.
\begin{itemize}
\item Classically, the cost is $O(n)=O(\mathrm{tol}^{-d/2})$ with optimal preconditioning.
\item On a quantum computer, Theorem~\ref{thm:main} gives $\mathrm{polylog}(1/\mathrm{tol})$
  calls per solve. Readout adds a factor $O(\mathrm{tol}^{-1})$, for a total of
  $\widetilde O(\mathrm{tol}^{-1})$.
\end{itemize}
An asymptotic advantage therefore requires $d=3$ on this count. This is
consistent with Deiml and Peterseim \cite[Remark~6.7]{DP26}. In their
analysis, if the accuracy demands $h\approx\mathrm{tol}^{1/s}$, an advantage requires
$d>s$. They measure the error in the energy norm, where $s=1$, which gives
an advantage from $d=2$. For quantities of interest with $O(h^2)$ accuracy,
$s=2$, so $d=3$ is needed. \todo{check the precise statement of their
Remark~6.7.}

\paragraph{Heat control.}
\begin{itemize}
\item With backward Euler, $N_t\sim\mathrm{tol}^{-1}$ and the classical cost is
  $\mathrm{tol}^{-(d/2+1)}$. On that count, a quantum advantage would appear for
  every $d$.
\item This comparison is unfair. For data that are smooth in time, a
  multiharmonic discretization needs only $\log(1/\mathrm{tol})$ harmonics, and the
  elliptic count returns.
\item The time dimension favours the quantum solver only when many
  harmonics are genuinely needed, i.e. for targets or forcing that are rough
  in time. \todo{quantify.}
\end{itemize}

\paragraph{Limitations.}
\begin{itemize}
\item The algorithms need fault-tolerant hardware.
\item The constants ($C_d$, the mass-matrix floor $3^d/(1-4^{-s})$, readout)
  matter in practice.
\item State preparation for general data and the end-to-end success
  probability need a full analysis.
\item The robustness of the multilevel factor in $\beta$ is proven
  (Theorem~\ref{thm:robust}) only for $Q_1$ elements on uniform grids, and
  with pessimistic constants.
\end{itemize}

\section{Conclusions and outlook}\label{sec:conclusions}

Preconditioning for quantum linear solvers is a numerical linear algebra
problem with its own rules. Respecting them leads to factorized,
well-scaled preconditioned operators whose effective condition numbers
are robust in the discretization parameters. Open directions:
\begin{enumerate}
\item a proof of $\beta$-robust multilevel bounds, and a better mass factor;
\item removing the $\sqrt L$ factor and computing explicit constants;
\item control constraints, where semismooth Newton or interior point
  methods change the $(2,2)$ block in every iteration. This links to
  preconditioners for interior point methods and to quantum interior point
  methods \cite{QIPMreview};
\item initial value problems. There $D_t$ is not circulant. The
  $\alpha$-circulant preconditioner \cite{MPW18} needs the non-unitary scaling
  $\Gamma_\alpha=\operatorname{diag}(\alpha^{k/N_t})$, whose condition number $1/\alpha$ enters the
  subnormalization. This trade-off between preconditioner quality and
  encoding cost has no classical analogue. Wavelet-in-time bases
  \cite{SS09} are an alternative;
\item variable coefficients, boundary control, and non-Cartesian meshes.
\end{enumerate}

\begin{thebibliography}{99}
\bibitem{BHMT02} G.~Brassard, P.~H{\o}yer, M.~Mosca, A.~Tapp, Quantum amplitude
  amplification and estimation, \emph{Contemp. Math.} 305 (2002) 53--74.
\bibitem{BNWA23} M.~Bagherimehrab, K.~Nakaji, N.~Wiebe, A.~Aspuru-Guzik, Fast
  quantum algorithm for differential equations, arXiv:2306.11802 (2023).
\bibitem{BHJ21} C.~Bacuta, D.~Hayes, J.~Jacavage, Efficient discretization and
  preconditioning of the singularly perturbed reaction-diffusion problem,
  arXiv:2104.11337 (2021).
\bibitem{BPX90} J.~H.~Bramble, J.~E.~Pasciak, J.~Xu, Parallel multilevel
  preconditioners, \emph{Math. Comp.} 55 (1990) 1--22.
\bibitem{BPV00} J.~H.~Bramble, J.~E.~Pasciak, P.~S.~Vassilevski, Computational
  scales of Sobolev norms with application to preconditioning, \emph{Math.\
  Comp.} 69 (2000) 463--480.
\bibitem{CJS13} B.~D.~Clader, B.~C.~Jacobs, C.~R.~Sprouse, Preconditioned quantum
  linear system algorithm, \emph{Phys. Rev. Lett.} 110 (2013) 250504.
\bibitem{CKS17} A.~M.~Childs, R.~Kothari, R.~D.~Somma, Quantum algorithm for
  systems of linear equations with exponentially improved dependence on
  precision, \emph{SIAM J. Comput.} 46 (2017) 1920--1950.
\bibitem{Costa22} P.~C.~S.~Costa, D.~An, Y.~R.~Sanders, Y.~Su, R.~Babbush,
  D.~W.~Berry, Optimal scaling quantum linear-systems solver via discrete
  adiabatic theorem, \emph{PRX Quantum} 3 (2022) 040303.
\bibitem{DMWL21} Y.~Dong, X.~Meng, K.~B.~Whaley, L.~Lin, Efficient phase-factor
  evaluation in quantum signal processing, \emph{Phys. Rev. A} 103 (2021)
  042419.
\bibitem{DP26} M.~Deiml, D.~Peterseim, Quantum realization of the finite element
  method, \emph{Math. Comp.} 95 (2026), no.~362, 2637--2668,
  doi:10.1090/mcom/4137; arXiv:2403.19512.
\bibitem{GO95} M.~Griebel, P.~Oswald, Tensor product type subspace splittings
  and multilevel iterative methods for anisotropic problems, \emph{Adv.\
  Comput.\ Math.} 4 (1995) 171--206.
\bibitem{GO95b} M.~Griebel, P.~Oswald, On the abstract theory of additive and
  multiplicative Schwarz algorithms, \emph{Numer.\ Math.} 70 (1995) 163--180.
\bibitem{GSLW19} A.~Gily\'en, Y.~Su, G.~H.~Low, N.~Wiebe, Quantum singular value
  transformation and beyond, \emph{Proc. STOC} (2019) 193--204.
\bibitem{HJ13} R.~A.~Horn, C.~R.~Johnson, \emph{Matrix Analysis}, 2nd ed.,
  Cambridge University Press, 2013.
\bibitem{JLMY25} S.~Jin, N.~Liu, C.~Ma, Y.~Yu, Quantum preconditioning method for
  finite difference discretizations of the Poisson equation via
  Schr\"odingerization, arXiv:2505.06866 (2025).
\bibitem{KXZ16} T.~Kolev, J.~Xu, Y.~Zhu, Multilevel preconditioners for
  reaction-diffusion problems with discontinuous coefficients, \emph{J.\
  Sci.\ Comput.} 67 (2016) 324--350.
\bibitem{KK12} M.~Kollmann, M.~Kolmbauer, A preconditioned MinRes solver for
  time-periodic parabolic optimal control problems, \emph{Numer. Linear
  Algebra Appl.} (2012), doi:10.1002/nla.1842.
\bibitem{KKLWZ13} M.~Kollmann, M.~Kolmbauer, U.~Langer, M.~Wolfmayr,
  W.~Zulehner, A robust finite element solver for a multiharmonic parabolic
  optimal control problem, \emph{Comput. Math. Appl.} 65 (2013) 469--486.
\bibitem{Lin22} L.~Lin, Lecture notes on quantum algorithms for scientific
  computation, arXiv:2201.08309 (2022).
\bibitem{LMS22} N.~Linden, A.~Montanaro, C.~Shao, Quantum vs.\ classical
  algorithms for solving the heat equation, \emph{Comm. Math. Phys.} 395
  (2022) 601--641.
\bibitem{LS25} L.~Lapworth, C.~S\"underhauf, Preconditioned block encodings for
  quantum linear systems, arXiv:2502.20908 (2025).
\bibitem{MP16} A.~Montanaro, S.~Pallister, Quantum algorithms and the finite
  element method, \emph{Phys. Rev. A} 93 (2016) 032324.
\bibitem{MPW18} E.~McDonald, J.~Pestana, A.~Wathen, Preconditioning and
  iterative solution of all-at-once systems for evolutionary partial
  differential equations, \emph{SIAM J. Sci. Comput.} 40 (2018)
  A1012--A1033.
\bibitem{Oswald94} P.~Oswald, \emph{Multilevel Finite Element Approximation:
  Theory and Applications}, Teubner, 1994.
\bibitem{PW12} J.~W.~Pearson, A.~J.~Wathen, A new approximation of the Schur
  complement in preconditioners for PDE-constrained optimization,
  \emph{Numer. Linear Algebra Appl.} 19 (2012) 816--829.
\bibitem{QIPMreview} M.~Mohammadisiahroudi et al., Quantum interior point
  methods: a review of developments and an optimally scaling framework,
  arXiv:2512.06224 (2025).
\bibitem{RDW10} T.~Rees, H.~S.~Dollar, A.~J.~Wathen, Optimal solvers for
  PDE-constrained optimization, \emph{SIAM J. Sci. Comput.} 32 (2010)
  271--298.
\bibitem{SS09} C.~Schwab, R.~Stevenson, Space-time adaptive wavelet methods for
  parabolic evolution problems, \emph{Math. Comp.} 78 (2009) 1293--1318.
\bibitem{SZ07} J.~Sch\"oberl, W.~Zulehner, Symmetric indefinite preconditioners
  for saddle point problems with applications to PDE-constrained
  optimization problems, \emph{SIAM J. Matrix Anal. Appl.} 29 (2007)
  752--773.
\bibitem{SuppProof} M.~Stoll, J.~W.~Pearson, Robustness of the weighted BPX
  factor for $M+\sqrt\beta K$, supplementary note (2026).
\bibitem{TAWL21} Y.~Tong, D.~An, N.~Wiebe, L.~Lin, Fast inversion, preconditioned
  quantum linear system solvers, fast Green's-function computation, and
  fast evaluation of matrix functions, \emph{Phys. Rev. A} 104 (2021)
  032422.
\bibitem{Xu92} J.~Xu, Iterative methods by space decomposition and subspace
  correction, \emph{SIAM Rev.} 34 (1992) 581--613.
\bibitem{XZ02} J.~Xu, L.~Zikatanov, The method of alternating projections and
  the method of subspace corrections in Hilbert space, \emph{J.\ Amer.\
  Math.\ Soc.} 15 (2002) 573--597.
\bibitem{YYZ26} Y.~Yang, Y.~Yu, L.~Zhang, A structure-aware construction of quantum
  BPX preconditioners for $Q_1$ finite element discretizations of
  reaction--diffusion problems, arXiv:2610.06987 (2026).
\end{thebibliography}

\end{document}
